\documentclass[UTF8,12pt]{ctexart}
\usepackage[a4paper,margin=24mm]{geometry}
\usepackage{amsmath,amssymb,bm,graphicx,adjustbox}
\newcommand{\fitmath}[1]{\begin{adjustbox}{max width=\linewidth}$\displaystyle #1$\end{adjustbox}}
\setlength{\parindent}{0pt}
\setlength{\parskip}{.7em}
\sloppy
\allowdisplaybreaks
\begin{document}
\section*{1995 年考研数学三 · 参考答案}
\subsection*{1995 年真题参考答案（试卷 IV）}

\subsection*{一、填空题}

（1）\(\displaystyle (-1)^n\dfrac{\displaystyle 2n!}{\displaystyle (1+x)^{n+1}}\)。 （2）\(\displaystyle 2z\)。 （3）\(\displaystyle x+e^x+C\)，其中 \(\displaystyle C\) 为任意常数。


（4）\(\displaystyle \begin{pmatrix}\dfrac{\displaystyle 1}{\displaystyle 10}&0&0\\\dfrac{\displaystyle 1}{\displaystyle 5}&\dfrac{\displaystyle 1}{\displaystyle 5}&0\\\dfrac{\displaystyle 3}{\displaystyle 10}&\dfrac{\displaystyle 2}{\displaystyle 5}&\dfrac{\displaystyle 1}{\displaystyle 2}\end{pmatrix}\)。 （5）\(\displaystyle \dfrac{\displaystyle \bar X}{\displaystyle Q}\sqrt{n(n-1)}\)。


\subsection*{二、选择题}

（1）D。 （2）A。 （3）C。 （4）D。 （5）C。


三、\(\displaystyle f(x)\) 在 \(\displaystyle x=0\) 处连续且可导，且 \(\displaystyle f^{\prime}(0)=0\)。


四、\(\displaystyle f(x)=3e^{3x}-2e^{2x}\)。


五、\(\displaystyle \ln(1-x-2x^2)=\displaystyle \sum_{n=1}^{\infty}\dfrac{\displaystyle (-1)^{n+1}-2^n}{\displaystyle n}x^n\)，其收敛区间为 \(\displaystyle (-\dfrac{\displaystyle 1}{\displaystyle 2},\allowbreak \dfrac{\displaystyle 1}{\displaystyle 2})\)。


六、\(\displaystyle I=-\sqrt{\dfrac{\displaystyle \pi}{\displaystyle 2}}\)。


七、\(\displaystyle P_0=\dfrac{\displaystyle ab}{\displaystyle b-1},\allowbreak \ Q_0=\dfrac{\displaystyle c}{\displaystyle 1-b}\)。


八、（1）证明略。（2）\(\displaystyle \dfrac{\displaystyle \pi}{\displaystyle 2}\)。


九、证明略。（由 \(\displaystyle R(\mathrm I)=R(\mathrm{II})=3\) 可得，存在 \(\displaystyle k_1,\allowbreak k_2,\allowbreak k_3\)，使得 \(\displaystyle \alpha_4=k_1\alpha_1+k_2\alpha_2+k_3\alpha_3\)。）


十、（1）\(\displaystyle f(x_1,\allowbreak x_2,\allowbreak x_3)=(x_1,\allowbreak x_2,\allowbreak x_3)\begin{pmatrix}0&2&-2\\2&4&4\\-2&4&-3\end{pmatrix}(x_1,\allowbreak x_2,\allowbreak x_3)^{\mathrm T}\)。


（2）令 \(\displaystyle P=\begin{pmatrix}\dfrac{\displaystyle 2}{\displaystyle \sqrt5}&\dfrac{\displaystyle 1}{\displaystyle \sqrt{30}}&\dfrac{\displaystyle 1}{\displaystyle \sqrt6}\\0&\dfrac{\displaystyle 5}{\displaystyle \sqrt{30}}&-\dfrac{\displaystyle 1}{\displaystyle \sqrt6}\\-\dfrac{\displaystyle 1}{\displaystyle \sqrt5}&\dfrac{\displaystyle 2}{\displaystyle \sqrt{30}}&\dfrac{\displaystyle 2}{\displaystyle \sqrt6}\end{pmatrix}\)，\(\displaystyle P\) 为正交矩阵。在正交变换 \(\displaystyle x=Py\) 下，二次型 \(\displaystyle f\) 化为标准形 \(\displaystyle f=y_1^2+6y_2^2-6y_3^2\)。


\subsection*{试卷 IV（续）}

十一、（1）\(\displaystyle \alpha=0.94^n\)。（2）\(\displaystyle \beta=\dbinom{\displaystyle n}{\displaystyle 2} 0.94^{n-2}0.06^2\)。（3）\(\displaystyle \theta=1-0.94^n-n\cdot0.94^{n-1}\cdot0.06\)。


十二、\(\displaystyle F(x,\allowbreak y)=\begin{cases}0,\allowbreak &x<0\text{ 或 }y<0,\allowbreak \\x^2y^2,\allowbreak &0\leq x\leq1,\allowbreak \ 0\leq y\leq1,\allowbreak \\x^2,\allowbreak &0\leq x\leq1,\allowbreak \ y>1,\allowbreak \\y^2,\allowbreak &x>1,\allowbreak \ 0\leq y\leq1,\allowbreak \\1,\allowbreak &x>1,\allowbreak \ y>1.\end{cases}\)


\subsection*{1995 年真题参考答案（试卷 V）}

\subsection*{一、填空题}

（1）2。 （2）【同试卷 IV 第一、（2）题】 （3）【同试卷 IV 第一、（3）题】 （4）【同试卷 IV 第一、（4）题】 （5）\(\displaystyle \dfrac{\displaystyle 1}{\displaystyle 6}\)。


\subsection*{二、选择题}

（1）【同试卷 IV 第二、（1）题】 （2）【同试卷 IV 第二、（2）题】 （3）C。 （4）C。 （5）【同试卷 IV 第二、（5）题】


三、【同试卷 IV 第三题】


四、\(\displaystyle x(\arcsin x)^2+2\sqrt{1-x^2}\arcsin x-2x+C\)，其中 \(\displaystyle C\) 为任意常数。


五、【同试卷 IV 第八题】


六、【同试卷 IV 第七题】


七、证明略。（可考虑函数 F(x)=xf(x)，对 F(x) 使用拉格朗日中值定理。）


八、\(\displaystyle f(x,\allowbreak y)\) 在闭区域 \(\displaystyle D\) 上的极值为 \(\displaystyle f(2,\allowbreak 1)=4\)，最大值为 \(\displaystyle f(2,\allowbreak 1)=4\)，最小值为 \(\displaystyle f(4,\allowbreak 2)=-64\)。


九、当 \(\displaystyle \lambda\ne-2,\allowbreak 1\) 时，方程组有唯一解；当 \(\displaystyle \lambda=-2\) 时，方程组无解；当 \(\displaystyle \lambda=1\) 时，方程组有无穷多解，其全部解为 \(\displaystyle x=k_1(-1,\allowbreak 1,\allowbreak 0)^{\mathrm T}+k_2(-1,\allowbreak 0,\allowbreak 1)^{\mathrm T}+(-2,\allowbreak 0,\allowbreak 0)^{\mathrm T}\)，其中 \(\displaystyle k_1,\allowbreak k_2\) 为任意常数。


十、\(\displaystyle A=\begin{pmatrix}\dfrac{\displaystyle 7}{\displaystyle 3}&0&-\dfrac{\displaystyle 2}{\displaystyle 3}\\0&\dfrac{\displaystyle 5}{\displaystyle 3}&-\dfrac{\displaystyle 2}{\displaystyle 3}\\-\dfrac{\displaystyle 2}{\displaystyle 3}&-\dfrac{\displaystyle 2}{\displaystyle 3}&2\end{pmatrix}\)。


十一、【同试卷 IV 第十一题】


十二、证明略。（计算 Y 的分布函数 G(Y)。）

\end{document}
